% Regression: t2_renorm13 — imported current-behavior fixture from the handoff corpus.
% Formula: (a):  if the MPS transfer channel E is divisible, E = R . R,
% t2_renorm13 — RMP arXiv:2011.12127, fig. fig:renorm2D, file
% fig2_fig2-pepe_newfig13.pdf: the inverse renormalization
% (fine-graining) step for MPS and for PEPS.  Fresh attempt at the
% figure t2_finegrain.tex recorded as failing.
%
% Formula (a):  if the MPS transfer channel E is divisible, E = R . R,
% the site tensor splits,
%     A^s_{alpha beta}  -->  sum_gamma B^{s_1}_{alpha gamma}
%                                      C^{s_2}_{gamma beta},
% one 3-leg site becoming two 3-leg sites of the SAME bond dimension.
%
% Formula (b):  the 2D analogue would need
%     A^s_{l r u d}  -->  (2 x 2 patch of tensors of bond dimension
%                          sqrt(D)),
% one 5-leg site becoming a 2x2 sheared patch — impossible for integer
% D, which is the caption's point.
%
% Declarative reading:
%   (a) two one-row grids, physical=up, default dot skin; the map arrow
%       between pictures is algebra, not tensor ink.
%   (b) two tenkzlattice pictures, frame=oblique (the sheared sheet),
%       boundary legs (the cut virtual bonds dangle), physical=up;
%       sizes 1x1 --> 2x2.
\documentclass[varwidth=19cm, border=6pt]{standalone}
\usepackage{amsmath, amssymb}
\usepackage{tenkz}
\begin{document}
\[
  (a)\quad
  % one MPS site: 3-leg dot
  \begin{tenkz}[physical=up]
    \tn{}
  \end{tenkz}
  \;\longrightarrow\;
  % the split: two sites of the same bond dimension
  \begin{tenkz}[physical=up]
    \tn{} & \tn{}
  \end{tenkz}
  \qquad\qquad
  (b)\quad
  % one PEPS site: 5-leg tensor on the sheared sheet
  \begin{tenkzlattice}[rows=1, cols=1, boundary legs, physical=up,
                       frame=oblique]
  \end{tenkzlattice}
  \;\longrightarrow\;
  % the would-be sqrt(D) fine-graining: a 2x2 patch
  \begin{tenkzlattice}[rows=2, cols=2, boundary legs, physical=up,
                       frame=oblique]
  \end{tenkzlattice}
\]
\end{document}
